\chapter{TỔNG QUAN BÀI TOÁN VÀ TIẾP CẬN TRÍ TUỆ NHÂN TẠO}

\section{Khảo sát thực trạng quản lý kho hàng}
Trong chuỗi cung ứng hiện đại, kho bãi là trung tâm lưu chuyển dòng vật chất và vốn của doanh nghiệp. Qua quá trình khảo sát tại các đơn vị phân phối linh kiện phần cứng máy tính và hàng tiêu dùng, các doanh nghiệp vừa và nhỏ (SMEs) thường gặp phải các vấn đề nghiêm trọng:
\begin{itemize}
    \item \textbf{Độ trễ thông tin:} Dữ liệu nhập xuất ghi chép bằng sổ sách hoặc bảng tính Excel phân tán, dẫn đến báo cáo chậm từ 1 đến 3 ngày.
    \item \textbf{Nguy cơ sai số và cháy hàng ngầm (Stock-out):} Thủ kho không thể theo dõi liên tục hàng trăm mã hàng (SKU) để biết sản phẩm nào đang cạn kiệt.
    \item \textbf{Đọng vốn lưu động (Dead-stock):} Hàng hóa kém luân chuyển nằm nhiều tháng không được cảnh báo để thanh lý kịp thời.
\end{itemize}

\section{Định vị vai trò của AI trong hệ thống}
Nhóm 11 xây dựng giải pháp AI với các cam kết kỹ thuật thực tế:
\begin{enumerate}
    \item \textbf{Chống ảo giác (Anti-Hallucination):} Mọi nhận định, bảng biểu từ AI đều được neo vào dữ liệu thực tế trích xuất trực tiếp từ cơ sở dữ liệu quan hệ (Ground-Truth Context Injection).
    \item \textbf{Cơ chế dự phòng kép (Dual-layer Resilience):} Kết hợp linh hoạt giữa Google Gemini Cloud LLM và bộ máy tính toán cục bộ (Local Rule Engine) hoạt động độc lập ngay cả khi mất mạng.
\end{enumerate}
